In the first example the array consists of $5$ elements, initially equal to zero. Consider the state of the array after each operation:

\begin{itemize}
\item after the first operation ($L_1 = 1$, $R_1 = 3$, $D_1 = 2$): $2\ 2\ 2\ 0\ 0$
\item after the second operation ($L_2 = 2$, $R_2 = 5$, $D_2 = -1$): $2\ 1\ 1\ -1\ -1$
\item after the third operation ($L_3 = 4$, $R_3 = 4$, $D_3 = 5$): $2\ 1\ 1\ 4\ -1$
\end{itemize}

The final array after applying all operations: ``2 1 1 4 -1''.

In the second example the array consists of a single element equal to zero. The only operation adds $7$ to the only element of the array ($L_1 = R_1 = 1$, $D_1 = 7$), so the final answer is ``7''.
